\scanpage{42}

d'après l'hypothèse \deleted{les $\unsure{s|\varphi|f}^{-1}(U_i)$ \uline{recouvrent} $T$ ;} \deleted{pour tout} $s\in S$, il existe un $i(s)$ tel que $\varphi(T_s)\subset U_{i(s)}$ et il en résulte (préciser) que l'on peut trouver un voisinage \deleted{\unsure{ou|un}} $V_{(s)}$ de $s$ tel que $\varphi(T\mid V_{(s)})\subset U_i$. On en déduit un hom $V(s)\longrightarrow H_i$, d'où $V(s)\unsure{-|\longrightarrow} H$. Il faut montrer que ça se recolle. C'est straight-forward \unsure{\ldots}.

\uline{Théorème} On suppose $T$ fini et loc. libre sur $S$, \deleted{et $X/S$ quelconque. Pour que $\operatorname{Hom}$} On suppose de plus que \deleted{\unsure{$T/S$ soit de|$T\to S$ soit de} l'\uline{une} des deux conditions suivantes (i) $T\unsure{-|\to} S$ est géom. injectif (ii)}

\margin{$k(s)\qquad k(t)$}

Pour tout $t\in S$, tte multisection de $X_s$ sur $T_s$, définie sur une extension finie de $k(s)$, soit contenue dans un ouvert affine de $X$. Alors \deleted{$\operatorname{Hom}$} \deleted{[condition tjrs remplie si tte partie} $f_*(X/T)$ existe, et est réunion des $f_{i*}(U_i/T)$, où $U_i$ parcourt les ouverts affines de $X$.

% archive sheet number: 41
